\documentclass[11pt]{article}

\usepackage[T1]{fontenc}
\usepackage[utf8]{inputenc}
\usepackage{lmodern}
\usepackage[final]{microtype}
\usepackage{geometry}
\geometry{margin=1in}
\usepackage{amsmath,amssymb,mathtools}
\usepackage{mathrsfs}
\usepackage{amsthm}
\usepackage{booktabs,tabularx,longtable,array}
\usepackage{enumitem}
\usepackage{xcolor}
\usepackage{hyperref}
\usepackage{xurl}
\usepackage{bookmark}
\usepackage{fancyhdr}
\usepackage{setspace}

\hypersetup{
  colorlinks=true,
  linkcolor=blue!45!black,
  citecolor=green!35!black,
  urlcolor=blue!55!black,
  pdftitle={Final resolution of the Fourier-decay problem for Rvachev's up-function},
  pdfauthor={OpenAI, prepared for Vladimir Reshetnikov}
}

\pagestyle{fancy}
\fancyhf{}
\fancyhead[L]{Fourier decay of Rvachev's up-function}
\fancyhead[R]{Final audit and resolution}
\fancyfoot[C]{\thepage}
\setlength{\headheight}{14pt}
\setlength{\emergencystretch}{3em}
\allowdisplaybreaks
\setlist[itemize]{leftmargin=2em,itemsep=0.25em,topsep=0.4em}
\setlist[enumerate]{leftmargin=2.25em,itemsep=0.3em,topsep=0.4em}

% Canonical project-wide mathematical notation.
% Canonical mathematical notation for active FabiusFunction LaTeX documents.
%
% This file is intentionally a plain TeX input rather than a LaTeX package.
% Every standalone document loads it by an explicit path relative to that
% document's directory; included fragments inherit the definitions from their
% root document.  Keep this file independent of document layout and metadata.
%
% Contract:
%   * one command has one mathematical meaning throughout the active corpus;
%   * names encode semantic roles rather than a document-local choice of letter;
%   * visually similar but differently normalized objects have different print;
%   * every command below is defined normatively in the notation catalogue.
%
% The active corpus excludes docs/archive/.  Historical sources may retain
% their original notation only inside an explicitly labelled quotation.

\makeatletter
\@ifundefined{mathbb}{%
  \PackageError{fabius-notation}{The amssymb package must be loaded first}%
    {Load amsmath and amssymb before inputting fabius-notation.tex.}%
}{}
\@ifundefined{operatorname}{%
  \PackageError{fabius-notation}{The amsmath package must be loaded first}%
    {Load amsmath before inputting fabius-notation.tex.}%
}{}
\makeatother

% Version stamp used by the catalogue and by static audits.
\newcommand{\FabiusNotationVersion}{2026-08-31}

% Number systems.  NaturalNumbers is explicitly the nonnegative convention.
\newcommand{\NaturalNumbers}{\mathbb{N}_{0}}
\newcommand{\PositiveIntegers}{\mathbb{N}_{>0}}
\newcommand{\IntegerNumbers}{\mathbb{Z}}
\newcommand{\RationalNumbers}{\mathbb{Q}}
\newcommand{\RealNumbers}{\mathbb{R}}
\newcommand{\ComplexNumbers}{\mathbb{C}}
\newcommand{\ExtendedRealNumbers}{\overline{\mathbb{R}}}

% The three principal Fabius--Rvachev functions.  Subscripts are deliberate:
% a bare F and a calligraphic F have several unrelated uses in the corpus.
\newcommand{\FabiusBounded}{F_{\mathrm{bd}}}
\newcommand{\FabiusGlobal}{F_{\mathrm{gl}}}
\newcommand{\FabiusQuantile}{Q_{F}}
\newcommand{\FabiusClampedQuantile}{Q_{F,\mathrm{cl}}}
\newcommand{\RvachevUp}{\operatorname{up}}

% Constants, elementary operators, and delimiters.
\newcommand{\EulerE}{\mathrm{e}}
\newcommand{\ImaginaryUnit}{\mathrm{i}}
\newcommand{\Differential}{\mathop{}\!\mathrm{d}}
\newcommand{\AbsoluteValue}[1]{\left\lvert #1\right\rvert}
\newcommand{\Norm}[1]{\left\lVert #1\right\rVert}
\newcommand{\Floor}[1]{\left\lfloor #1\right\rfloor}
\newcommand{\Ceiling}[1]{\left\lceil #1\right\rceil}
\newcommand{\FractionalPart}[1]{\left\{#1\right\}}
\newcommand{\PositivePart}[1]{\left(#1\right)_{+}}
\newcommand{\IndicatorOf}[1]{\mathbf{1}_{#1}}
\newcommand{\SetOf}[1]{\left\{#1\right\}}
\newcommand{\InnerProduct}[1]{\left\langle #1\right\rangle}
\newcommand{\CardinalityOf}[1]{\left\lvert #1\right\rvert}
\newcommand{\DefinitionEquals}{\mathrel{\coloneqq}}
\newcommand{\Sign}{\operatorname{sgn}}
\newcommand{\IdentityOperator}{\operatorname{id}}
\newcommand{\SpanOperator}{\operatorname{span}}
\newcommand{\TraceOperator}{\operatorname{tr}}
\newcommand{\RankOperator}{\operatorname{rank}}
\newcommand{\DiagonalOperator}{\operatorname{diag}}
\newcommand{\ResidueOperator}{\operatorname*{Res}}
\newcommand{\LeastCommonMultiple}{\operatorname{lcm}}
\newcommand{\OrderOperator}{\operatorname{ord}}
\newcommand{\PrincipalLogarithm}{\operatorname{Log}}
\newcommand{\FallingFactorial}[2]{\left(#1\right)^{\underline{#2}}}
\newcommand{\RisingFactorial}[2]{\left(#1\right)^{\overline{#2}}}

% Support, topology, convolution, and metric notation.
\newcommand{\SupportOperator}{\operatorname{supp}}
\newcommand{\SupportOf}[1]{\SupportOperator\!\left(#1\right)}
\newcommand{\TopologicalSupportOf}[1]{\operatorname{tsupp}\!\left(#1\right)}
\newcommand{\Convolution}{\mathbin{*}}
\newcommand{\MetricDistance}{\operatorname{dist}}
\newcommand{\TotalVariationDistance}{d_{\mathrm{TV}}}

% Probability.  Equality and convergence in law are not metric distance.
\newcommand{\Probability}{\mathbb{P}}
\newcommand{\Expectation}{\mathbb{E}}
\newcommand{\Variance}{\operatorname{Var}}
\newcommand{\Covariance}{\operatorname{Cov}}
\newcommand{\LawOperator}{\operatorname{Law}}
\newcommand{\LawOf}[1]{\LawOperator\!\left(#1\right)}
\newcommand{\EqualInLaw}{\mathrel{\overset{\mathrm{law}}{=}}}
\newcommand{\ConvergesInLaw}{\mathrel{\xrightarrow{\mathrm{law}}}}

% Fourier and integral transforms.  The printed labels expose normalization.
% FourierTwoPi uses exp(-2*pi*i*x*xi); FourierAngular uses exp(-i*x*omega).
\newcommand{\FourierTwoPi}[1]{\widehat{#1}_{\,2\pi}}
\newcommand{\FourierAngular}[1]{\widehat{#1}_{\,\mathrm{ang}}}
\newcommand{\LaplaceTransformSymbol}{\mathcal{L}}
\newcommand{\MellinTransformSymbol}{\mathcal{M}}
\newcommand{\LaplaceTransformOf}[1]{\LaplaceTransformSymbol\!\left[#1\right]}
\newcommand{\MellinTransformOf}[1]{\MellinTransformSymbol\!\left[#1\right]}
\newcommand{\MomentGeneratingFunction}[1]{M_{#1}}
\newcommand{\CharacteristicFunction}[1]{\varphi_{#1}}

% The two standard sinc functions must never share one symbol.
\newcommand{\SincRad}{\operatorname{sinc}_{\mathrm{rad}}}
\newcommand{\SincPi}{\operatorname{sinc}_{\pi}}

% Binary arithmetic and Thue--Morse notation.
\newcommand{\BinaryDigitSum}{\operatorname{wt}_{2}}
\newcommand{\ThueMorseSignSymbol}{\varepsilon}
\newcommand{\ThueMorseSign}[1]{\varepsilon_{#1}}
\newcommand{\PadicValuation}[1]{\nu_{#1}}
\newcommand{\TwoAdicValuation}{\nu_{2}}

% q-series notation.  Every base is an explicit argument.
\newcommand{\QPochhammer}[3]{\left(#1;#2\right)_{#3}}
\newcommand{\GaussianBinomial}[3]{%
  \genfrac{[}{]}{0pt}{}{#1}{#2}_{#3}%
}
\newcommand{\QInteger}[2]{\left[#1\right]_{#2}}

% Named special functions.
\newcommand{\Polylogarithm}{\operatorname{Li}}
\newcommand{\LambertW}{W}
\newcommand{\GammaFunction}{\Gamma}
\newcommand{\RiemannZeta}{\zeta}

% Asymptotics.  BigO and LittleO are operator tokens for existing formula
% grammar; prefer the At forms so that the limiting regime is printed.
\newcommand{\BigO}{\mathcal{O}}
\newcommand{\LittleO}{o}
\newcommand{\BigOAt}[2]{\mathcal{O}_{#2}\!\left(#1\right)}
\newcommand{\LittleOAt}[2]{o_{#2}\!\left(#1\right)}

\newtheorem{theorem}{Theorem}[section]
\newtheorem{proposition}[theorem]{Proposition}
\newtheorem{lemma}[theorem]{Lemma}
\newtheorem{corollary}[theorem]{Corollary}
\newtheorem{claim}[theorem]{Claim}
\theoremstyle{definition}
\newtheorem{definition}[theorem]{Definition}
\newtheorem{problem}[theorem]{Problem}
\theoremstyle{remark}
\newtheorem{remark}[theorem]{Remark}
\newtheorem{warning}[theorem]{Warning}

\newcommand{\Leb}{\operatorname{Leb}}
\newcommand{\T}{T}
\newcommand{\Ecal}{\mathcal E}
\newcommand{\Wlam}{W}
\newcommand{\ind}{\mathbf 1}
\newcommand{\Ceq}{\textup{(C)}}
\newcommand{\Beq}{\textup{(B)}}

\newenvironment{verdict}{%
  \par\medskip\noindent\begin{center}\begin{minipage}{0.94\textwidth}\hrule\vspace{0.65em}\small
}{%
  \vspace{0.65em}\hrule\end{minipage}\end{center}\medskip
}

\title{\bfseries Final Resolution of the Real-Axis Fourier-Decay Problem\\
for Rvachev's Up-Function\\[0.4em]
\large Lacunary spectra, a corrected iterated-logarithm law, sharp Lambert--\(W\) peaks,\\
and natural versus shell-adaptive Ces\`aro gauges}
\author{Prepared for Vladimir Reshetnikov's \texttt{ProveIt} project}
\date{26 August 2026}

\begin{document}
\maketitle

\begin{abstract}
Let
\[
 \Phi(x)=\widehat{\RvachevUp}(x)
 =\prod_{h=0}^{\infty}\SincRad\!\left(\frac{\pi x}{2^h}\right),
 \qquad \SincRad z=\frac{\sin z}{z},
\]
be the Fourier transform of Rvachev's compactly supported up-function under the
\(\EulerE^{-2\pi\ImaginaryUnit x\xi}\) convention.  This article audits the original question,
four independent first-wave analyses, the Aistleitner--Hofer--Larcher paper on
lacunary trigonometric products, and the subsequent comparative audit.  The
apparently conflicting answers are reconciled by an exact dyadic factorization:
a universal triangular denominator produces
\(\exp[-(\log x)^2/(2\log2)]\), while a Thue--Morse sine cocycle supplies a
whole spectrum of power corrections.  The typical, arithmetic-mean,
root-mean-square, and uniform exponents are respectively
\[
 \begin{gathered}
 \kappa_0=3.1514961294\ldots,\qquad
 \kappa_1=2.7480700149\ldots,\\
 \kappa_2=2.6514961294\ldots,\qquad
 \kappa_\infty=2.3590148791\ldots.
 \end{gathered}
\]
We give a direct covariance and martingale calculation showing that the correct
per-level variance of \(\log|2\sin\pi x|\) under doubling is \(\pi^2/4\), not
\(\pi^2/2\).  Thus the law-of-the-iterated-logarithm constants printed in the
2017 lacunary-products paper, and inherited by one first-wave report, must be
divided by \(\sqrt2\).

For the original arithmetic Ces\`aro question, the natural gauge
\[
 g_0(x)=A_1x^{-\kappa_1}
 \exp\!\left(-\frac{(\log x)^2}{2\log2}\right)
\]
solves the requested two-sided boundedness and gives limit one at dyadic
endpoints, but its unrestricted averages have a nonconstant dyadic-mantissa
profile governed by a continuous singular eigenmeasure.  This proves
nonexistence of an unrestricted limit for every shell-shape-stabilizing gauge,
in particular for the usual Hardy-field or logarithmico-exponential closed
forms.  The comparative audit stopped there and consequently overstated its
scope.  We close the literal existential problem: by progressively flattening
the shell measures on finer finite partitions, while using the very steep
baseline slope to preserve monotonicity, one obtains a positive \(C^\infty\)
decreasing gauge with the all-\(t\) limit.  A Carleman--Hoischen approximation
of the reciprocal and its derivative upgrades it to a positive real-analytic,
strictly decreasing gauge.  It has the same rough asymptotic
\(g(x)=g_0(x)x^{o(1)}\), but its within-shell shape necessarily becomes more
complex and does not stabilize.

Finally, for the actual largest value
\(M_k=\max_{2^k\le x\le2^{k+1}}|\Phi(x)|\), we verify the sharp boundary-layer
formula
\[
 \begin{aligned}
 \log M_k={}&-\frac{\log2}{2}k(k+1)-(k+1)\log\pi
 +(k+1)\log\frac{\sqrt3}{2}\\
 &-\frac{W(c(k+1))^2+2W(c(k+1))}{2\log2}
 +\BigO((\log\log k)^2),
 \end{aligned}
\]
where \(c=(\log2)\sqrt{3/2}/\pi\).  This includes a second log-normal
suppression in \(\log\log x\), and is distinct from every fixed-ray or mean
law.
\end{abstract}

\begingroup
\footnotesize
\tableofcontents
\endgroup
\newpage

\section{The question and the final verdict}

Set
\begin{equation}
 \Phi(x)=\prod_{h=0}^{\infty}
 \SincRad\!\left(\frac{\pi x}{2^h}\right),
 \qquad x\ge0,
 \label{eq:Phi-def}
\end{equation}
and let \(t_0>0\) be fixed.  The original question asks for a real-valued
function \(g\), preferably elementary, real-analytic, positive and eventually
decreasing, for which\cite{mse-question,question-source}
\begin{equation*}
 \frac1{t-t_0}\int_{t_0}^{t}\frac{|\Phi(x)|}{g(x)}\Differential x\longrightarrow1
 \qquad(t\to\infty).
 \tag{C}
\end{equation*}
or at least for which the same expression stays between two positive constants:
\begin{equation*}
 0<A<\frac1{t-t_0}\int_{t_0}^{t}\frac{|\Phi(x)|}{g(x)}\Differential x<B<\infty
 \qquad(t\ge t_1).
 \tag{B}
\end{equation*}
The adjective ``decay rate'' is ambiguous unless one specifies whether the
quantity being measured is pointwise, typical, averaged in an \(L^q\) norm, or
maximized over a dyadic annulus.  The exact answer is therefore a dictionary,
not a single equivalent.

Write
\begin{equation}
 a:=\log2,
 \qquad
 \Ecal_\kappa(x):=x^{-\kappa}
 \exp\!\left(-\frac{(\log x)^2}{2a}\right).
 \label{eq:E-kappa}
\end{equation}
The four principal exponents are
\begin{align}
 \kappa_0&=\frac32+\frac{\log\pi}{a}
          =3.1514961294723188\ldots, \label{eq:kappa0}\\
 \kappa_1&=\frac12+\frac{\log(\pi/\rho_1)}{a}
          =2.7480700148713340\ldots, \label{eq:kappa1}\\
 \kappa_2&=\frac{\log(2\pi)}{a}
          =2.6514961294723188\ldots, \label{eq:kappa2}\\
 \kappa_\infty&=\frac32+\frac{\log(\pi/\sqrt3)}{a}
          =2.3590148791117407\ldots, \label{eq:kappainf}
\end{align}
where \(\rho_1\) is the Perron root of the weighted doubling operator in
\eqref{eq:L1}; numerically
\begin{equation}
 \rho_1=0.661322602060565\ldots .
 \label{eq:rho1-num}
\end{equation}
The decimal is a high-precision numerical value.  Rigorous enclosures include
\(0.66130<\rho_1<0.66135\) from Aistleitner--Hofer--Larcher and the independent
exact-arithmetic certificate
\(0.66126807<\rho_1<0.66134891\) in the fourth first-wave report\cite{ahl2017,wave1-4}.

\begin{theorem}[Final resolution]
\label{thm:final}
The following statements hold.
\begin{enumerate}[label=\textup{(\roman*)}]
\item There is no positive pointwise equivalent for \(|\Phi(x)|\) on the full
real ray, because every nonzero integer is a zero.  Nevertheless
\[
 \limsup_{x\to\infty}\frac{\log|\Phi(x)|}{(\log x)^2}
 =-\frac1{2\log2},
 \qquad
 \liminf_{x\to\infty}\frac{\log|\Phi(x)|}{(\log x)^2}=-\infty.
\]
\item On almost every fixed dyadic ray \(x=2^ky\), \(1<y<2\), the exponent is
\(\kappa_0\).  The centered logarithm satisfies a CLT with variance
\(\pi^2/4\) per dyadic level and the LIL constants in
\eqref{eq:correct-LIL}.  This corrects a factor \(\sqrt2\) in the cited
2017 paper.
\item The normalized \(L^q\) shell amplitudes form the pressure spectrum
\(\kappa_q\) of \eqref{eq:kappa-q}.  In particular, \(q=1\) gives the
arithmetic-mean exponent \(\kappa_1\), \(q=2\) gives exactly
\(\kappa_2=\log_2(2\pi)\), and \(q=\infty\) gives \(\kappa_\infty\).
\item The natural analytic gauge \(g_0=A_1\Ecal_{\kappa_1}\), with \(A_1>0\)
defined in \eqref{eq:A1}, satisfies \Beq; at \(t=2^k\) its
Ces\`aro ratio tends to one.  At \(t=2^ky\), the ratio tends to the
nonconstant profile \(\mathscr P(y)\) of \eqref{eq:profile}.  Consequently
\Ceq fails for every shell-shape-stabilizing gauge, including all
Hardy-field gauges of the correct rough size and their continuous log-periodic
modulations.
\item In the literal unrestricted existential sense, \Ceq is
\emph{affirmative}.  There exists a positive real-analytic strictly decreasing
function \(g_{\mathrm{ad}}\) on \((0,\infty)\) such that \Ceq holds and
\begin{equation}
 \log g_{\mathrm{ad}}(x)
 =-\frac{(\log x)^2}{2\log2}-\kappa_1\log x+o(\log x).
 \label{eq:adaptive-rough}
\end{equation}
Its correction relative to \(g_0\) is shell-adaptive and cannot have a
stabilizing within-shell profile.
\item The natural gauge becomes fully convergent without adaptation if the
arithmetic mean is replaced by the logarithmic mean \(\Differential x/x\), as in
\eqref{eq:logmean}.
\item The dyadic-annulus maximum obeys the Lambert--\(W\) asymptotic
\eqref{eq:annulus-main}.  Local peaks immediately following \(m2^k\), with
\(m\) fixed odd, obey the much smaller two-adic law
\eqref{eq:twoadic-peak}.
\end{enumerate}
\end{theorem}

\begin{verdict}
\textbf{What ``decays how fast'' means.}
The universal factor is
\(\exp[-(\log x)^2/(2\log2)]\).  Its power correction is
\(x^{-\kappa_0}\) at a typical dyadic phase,
\(x^{-\kappa_1}\) in absolute mean,
\(x^{-\kappa_2}\) in RMS, and
\(x^{-\kappa_\infty}\) for the sharp pointwise upper scale.
The largest peak in an entire dyadic annulus has an additional
Lambert--\(W\), equivalently \((\log\log x)^2\), boundary-layer loss.
For the original Ces\`aro question, a simple natural gauge gives boundedness
but a nonconstant phase profile; an exact all-time analytic gauge exists only
by allowing an increasingly nonstationary shell correction.
\end{verdict}

\section{Audit of the supplied materials}

The four first-wave documents are not mutually contradictory once their target
functionals are separated.  The substantive conflict is the LIL constant, and
the substantive gap in the comparative audit is its extension of a
class-restricted nonexistence theorem to the literal unrestricted question.

\begin{longtable}{@{}p{0.16\textwidth}p{0.35\textwidth}p{0.40\textwidth}@{}}
\toprule
Source & Principal contribution & Final assessment \\
\midrule
\endfirsthead
\toprule
Source & Principal contribution & Final assessment \\
\midrule
\endhead
First-wave 1 \cite{wave1-1} & Exact dyadic reduction, almost-everywhere exponent, CLT/LIL,
coarse and then refined annular behavior. & Its variance \(\pi^2/4\) and LIL
normalization are correct.  It does not solve the literal arithmetic Ces\`aro
question. \\
\addlinespace
First-wave 2 \cite{wave1-2} & Strongest elementary treatment of the \(L^\infty\) geometry:
calibrated Gelfond inequality, dense preimages of the maximizing cycle,
Lambert--\(W\) annular maximum, and two-adic valleys. & Correct and sharply
scoped.  Its warning that the original \(L^1\) question is separate is
essential. \\
\addlinespace
First-wave 3 \cite{wave1-3} & Arithmetic-mean exponent \(\kappa_1\), transfer operator,
bounded Ces\`aro gauge, dyadic profile, and the \(q\)-spectrum viewpoint. &
Correct on the mean theory, except that its LIL constant faithfully inherits
the erroneous printed constant from the 2017 paper. \\
\addlinespace
First-wave 4 \cite{wave1-4} & Most systematic \(L^q\) spectrum, exact \(L^2\) identity,
certificate for \(\rho_1\), singular eigenmeasure, unique lobe peaks, and
exceptional valleys. & Correct where checked; the RMS identification
\(\kappa_2=\log_2(2\pi)\) is exact.  Its natural-gauge result is not a proof of
nonexistence for arbitrary analytic gauges. \\
\addlinespace
Aistleitner--Hofer--Larcher \cite{ahl2017,ahl-source} & Published lacunary-product LIL and the
Fouvry--Mauduit constant \(\rho_1\). & The \(\rho_1\) results are unaffected.
Equation (4.5) omits the factor \(1/2\) in cosine Parseval, doubling the
variance.  The printed LIL constants must be divided by \(\sqrt2\). \\
\addlinespace
Comparative audit \cite{audit} & Correct synthesis of the natural gauge, singular phase
profile, class-restricted obstruction, logarithmic mean, and variance
correction. & The theorem actually excludes shell-stabilizing gauges (and
Hardy-field gauges), not every real-analytic decreasing gauge.  The statement
that the unrestricted limit is ``unattainable in principle'' is therefore too
strong.  Section~\ref{sec:adaptive} supplies a countervailing existence theorem. \\
\bottomrule
\end{longtable}

Two numerical cautions are worth recording.  First, digits such as
\(\rho_1=0.661322602060565\ldots\) and \(A_1=0.0912661\ldots\) arise from
spectral collocation and quadrature; the exact definitions, rather than the
decimals, are used in the theorems.  Second, the very slow approach to the
asymptotic regime near two-adic valleys makes small-scale numerical fitting a
poor guide to the final exponents.

\section{Exact dyadic factorization and lobe structure}
\label{sec:factorization}

The infinite product converges locally uniformly, because
\(\SincRad z=1-z^2/6+\BigO(z^4)\) at the origin.  Hence \(\Phi\) is even and
entire.  It also satisfies the exact refinement equation\cite{arias1982,wikipedia}
\begin{equation}
 \Phi(2x)=\SincRad(2\pi x)\Phi(x).
 \label{eq:refinement}
\end{equation}

For \(1\le y\le2\), put \(t=y-1\) and
\begin{align}
 P_n(t)&:=\prod_{j=0}^{n-1}|\sin(\pi2^jt)|, \label{eq:Pn}\\
 H(y)&:=\prod_{r=1}^{\infty}
 \left|\SincRad\!\left(\frac{\pi y}{2^r}\right)\right|.
 \label{eq:H}
\end{align}
The tail \(H\) is continuous on \([1,2]\), positive on \([1,2)\), and
\(H(2)=0\).

\begin{proposition}[Exact shell factorization]
\label{prop:factorization}
For every integer \(k\ge0\) and \(1\le y\le2\), with \(n=k+1\),
\begin{equation}
 |\Phi(2^ky)|
 =2^{-k(k+1)/2}(\pi y)^{-n}H(y)P_n(y-1).
 \label{eq:factorization}
\end{equation}
\end{proposition}

\begin{proof}
Split the factors in \eqref{eq:Phi-def} at level \(k\).  In the first
\(k+1\) factors, substitute \(r=k-h\):
\[
 \prod_{h=0}^{k}
 \left|\SincRad\!\left(\frac{\pi2^ky}{2^h}\right)\right|
 =\frac{\prod_{r=0}^{k}|\sin(\pi2^ry)|}
 {(\pi y)^{k+1}2^{0+1+\cdots+k}}.
\]
Since \(y=1+t\), the absolute numerator is
\(\prod_{r=0}^k|\sin(\pi2^rt)|=P_{k+1}(t)\).  The remaining factors are
exactly \(H(y)\).
\end{proof}

The factorization isolates the deterministic triangular denominator from a
multiplicative cocycle for the doubling map
\begin{equation}
 \T t=2t\pmod1,
 \qquad
 P_n(t)=\prod_{j=0}^{n-1}s(\T^jt),
 \qquad s(t)=|\sin\pi t|.
 \label{eq:doubling}
\end{equation}
It also gives the exact normalized identity
\begin{equation}
 \frac{|\Phi(2^ky)|}{\Ecal_\kappa(2^ky)}
 =C_\kappa(y)
 \exp\!\bigl(k(\kappa a-\log\pi-a/2)\bigr)P_{k+1}(y-1),
 \label{eq:master-normalization}
\end{equation}
where
\begin{equation}
 C_\kappa(y)=\frac{H(y)}{\pi}
 y^{\kappa-1}\exp\!\left(\frac{(\log y)^2}{2a}\right).
 \label{eq:Ckappa}
\end{equation}
All competing exponents arise by inserting a different exponential scale for
\(P_n\).

\subsection{Zeros, signs, and one extremum per lobe}

Euler's product for sine gives the canonical product
\begin{equation}
 \Phi(z)=\prod_{m=1}^{\infty}
 \left(1-\frac{z^2}{m^2}\right)^{1+\TwoAdicValuation(m)}.
 \label{eq:canonical-product}
\end{equation}
Thus a nonzero integer \(m\) is a zero of multiplicity
\(1+\TwoAdicValuation(|m|)\).  On every interval \((N,N+1)\),
\begin{equation}
 \frac{\Differential^2}{\Differential x^2}\log|\Phi(x)|
 =-2\sum_{m=1}^{\infty}(1+\TwoAdicValuation(m))
 \frac{m^2+x^2}{(m^2-x^2)^2}<0.
 \label{eq:strict-log-concavity}
\end{equation}
The first derivative tends to \(+\infty\) at the left endpoint and to
\(-\infty\) at the right endpoint.  Hence every lobe has exactly one critical
point, necessarily its strict maximum.  The sign on \((N,N+1)\) is
\((-1)^{s_2(N)}\), the signed Thue--Morse value; the present article concerns
its modulus.

\section{Typical fixed rays and the corrected fluctuation law}
\label{sec:typical}

Define the centered observable
\begin{equation}
 X(t):=\log|2\sin\pi t|,
 \qquad
 S_n(t):=\sum_{j=0}^{n-1}X(\T^jt).
 \label{eq:X-S}
\end{equation}
Since \(\int_0^1\log|\sin\pi t|\Differential t=-\log2\), one has
\(\int_0^1X=0\).  At \(\kappa=\kappa_0\),
\eqref{eq:master-normalization} becomes the exact formula
\begin{equation}
 |\Phi(2^ky)|
 =\Ecal_{\kappa_0}(2^ky)D_0(y)\EulerE^{S_{k+1}(y-1)},
 \qquad
 D_0(y)=\frac{H(y)}{2\pi}
 y^{\kappa_0-1}\EulerE^{(\log y)^2/(2a)}.
 \label{eq:typical-exact}
\end{equation}

\begin{theorem}[Almost-everywhere dyadic-ray law]
\label{thm:typical}
For Lebesgue-almost every \(y\in(1,2)\),
\begin{equation}
 \log|\Phi(2^ky)|
 =-\frac{(\log(2^ky))^2}{2a}
 -\kappa_0\log(2^ky)+o(k).
 \label{eq:typical-law}
\end{equation}
Moreover, with \(n=k+1\),
\begin{equation}
 \frac{S_n}{\sqrt n}\ \Longrightarrow\
 N\!\left(0,\frac{\pi^2}{4}\right),
 \label{eq:CLT}
\end{equation}
and almost surely
\begin{equation}
 \limsup_{n\to\infty}\frac{S_n}{\sqrt{2n\log\log n}}=\frac\pi2,
 \qquad
 \liminf_{n\to\infty}\frac{S_n}{\sqrt{2n\log\log n}}=-\frac\pi2.
 \label{eq:correct-LIL}
\end{equation}
Equivalently, in the normalization \(\sqrt{n\log\log n}\), the constants are
\(\pm\pi/\sqrt2\).
\end{theorem}

\begin{proof}
Birkhoff's theorem applied to the integrable potential
\(\log|\sin\pi t|\) gives \(S_n/n\to0\) almost everywhere, proving
\eqref{eq:typical-law}.  The exact variance and martingale reduction are proved
in the next two propositions; the standard stationary reverse-martingale CLT
and LIL\cite{hallheyde} then give \eqref{eq:CLT}--\eqref{eq:correct-LIL}.  The coboundary
remainder is negligible because \(X\) has exponentially decaying tails.
\end{proof}

\subsection{The exact covariance calculation}

In \(L^2(0,1)\),
\begin{equation}
 X(t)=-\sum_{m=1}^{\infty}\frac{\cos(2\pi mt)}{m}.
 \label{eq:Clausen}
\end{equation}
Cosine orthogonality gives, for \(r\ge0\),
\begin{equation}
 c_r:=\int_0^1X(t)X(2^rt)\Differential t
 =\frac{\pi^2}{12\,2^r}.
 \label{eq:covariance}
\end{equation}
Indeed, only the frequency pairs \(m=2^rj\) survive, and every surviving
cosine integral is \(1/2\), not \(1\).

\begin{proposition}[Exact finite variance]
\label{prop:variance}
For every \(n\ge1\),
\begin{equation}
 \Variance(S_n)
 =\frac{\pi^2}{4}n-\frac{\pi^2}{3}(1-2^{-n}).
 \label{eq:exact-variance}
\end{equation}
In particular, the asymptotic variance per level is \(\pi^2/4\).
\end{proposition}

\begin{proof}
Stationarity and \eqref{eq:covariance} yield
\[
 \Variance(S_n)=nc_0+2\sum_{r=1}^{n-1}(n-r)c_r.
\]
Using \(\sum_{r=1}^{n-1}(n-r)2^{-r}=n-2+2^{1-n}\) gives
\eqref{eq:exact-variance}.
\end{proof}

\subsection{A one-line martingale decomposition}

Let \(\mathcal P\) be the Perron operator of doubling,
\begin{equation}
 (\mathcal Ph)(x)=\frac12\left[h\!\left(\frac x2\right)
 +h\!\left(\frac{x+1}{2}\right)\right].
 \label{eq:Perron-unweighted}
\end{equation}
The identity
\begin{equation}
 \mathcal PX=\frac12X
 \label{eq:PX}
\end{equation}
follows from \(2\sin(\pi x/2)\,2\cos(\pi x/2)=2\sin\pi x\).
Set
\begin{equation}
 d:=2X-X\circ\T.
 \label{eq:d}
\end{equation}
Then \(\mathcal Pd=0\), and
\begin{equation}
 S_n=\sum_{j=0}^{n-1}d\circ\T^j+X\circ\T^n-X.
 \label{eq:martingale-decomp}
\end{equation}
The first term is a stationary reverse-martingale sum.  Furthermore
\begin{equation}
 \int_0^1d^2
 =4c_0+c_0-4c_1=\frac{\pi^2}{4}.
 \label{eq:d-variance}
\end{equation}
Since \(\Leb(|X|>u)=\BigO(\EulerE^{-u})\), Borel--Cantelli implies
\(X\circ\T^n=o(\sqrt{n\log\log n})\) almost surely.  This proves that the
martingale variance and LIL constants pass unchanged to \(S_n\).

\subsection{The factor-\texorpdfstring{\(\sqrt2\)}{sqrt(2)} error in the lacunary-products paper}

The cited Aistleitner--Hofer--Larcher source\cite{ahl2017,ahl-source} evaluates the same variance
functional as
\[
 \sum_{j\ge1}\left(\frac1{j^2}
 +2\sum_{r\ge1}\frac1{2^rj^2}\right)=\frac{\pi^2}{2}.
\]
That displayed sum is twice the correct Parseval expression: each pairing of
cosines contributes \(1/2\).  Therefore the source's equation (4.5) should be
\(\sigma_{f_1}^2=\pi^2/4\).  In the source's normalizations, its intermediate
LIL constant \(\pi/\sqrt2\) becomes \(\pi/2\), its theorem constant \(\pi\)
becomes \(\pi/\sqrt2\), and the corresponding product constant
\(\pi/\sqrt{\log2}\) becomes
\(\pi/\sqrt{2\log2}\).  The Fouvry--Mauduit constant and the paper's
\(L^1\)-based discrepancy exponent are unaffected.

\section{The full \texorpdfstring{\(L^q\)}{Lq} decay spectrum}
\label{sec:q-spectrum}

For \(q>0\), define the weighted transfer operator
\begin{equation}
 (\LaplaceTransformSymbol_qh)(x)=\frac12\sum_{\T y=x}s(y)^q h(y),
 \qquad s(y)=|\sin\pi y|.
 \label{eq:Lq}
\end{equation}
Thus
\begin{equation}
 \int_0^1P_n(t)^qh(t)\Differential t
 =\int_0^1\LaplaceTransformSymbol_q^nh(x)\Differential x.
 \label{eq:Lq-integral}
\end{equation}
Let \(\rho_q\) be the spectral radius of \(\LaplaceTransformSymbol_q\), and put
\begin{equation}
 \Lambda_q:=\rho_q^{1/q},
 \qquad
 \kappa_q:=\frac{\log\pi+a/2-\log\Lambda_q}{a}.
 \label{eq:kappa-q}
\end{equation}
Standard Ruelle--Perron--Frobenius theory on a suitable H\"older or
bounded-variation space\cite{baladi} gives positive eigenvectors and exponential
convergence for the endpoint-power weights occurring here.  It follows from
\eqref{eq:master-normalization} that the normalized \(L^q\) shell norm of
\(|\Phi|/\Ecal_{\kappa_q}\) tends to a finite positive constant.

The endpoint cases can be computed almost entirely by hand.

\subsection{Geometric mean and \texorpdfstring{\(q\downarrow0\)}{q to 0}}

Lebesgue measure is invariant under \(\T\), so
\begin{equation}
 \int_0^1\log P_n(t)\Differential t
 =n\int_0^1\log|\sin\pi t|\Differential t=-na.
 \label{eq:geometric-mean}
\end{equation}
Hence \(\Lambda_0=1/2\) and \(\kappa_0\) is \eqref{eq:kappa0}.  This is the
same exponent as the almost-everywhere fixed-ray law.

\subsection{Arithmetic mean and the Fouvry--Mauduit constant}

For \(q=1\), \eqref{eq:Lq} becomes
\begin{equation}
 (\LaplaceTransformSymbol_1h)(x)=\frac12\left[
 \sin\!\left(\frac{\pi x}{2}\right)h\!\left(\frac x2\right)
 +\cos\!\left(\frac{\pi x}{2}\right)h\!\left(\frac{x+1}{2}\right)
 \right].
 \label{eq:L1}
\end{equation}
Fouvry and Mauduit proved\cite{fouvrymauduit,ahl2017}
\begin{equation}
 I_1(n):=\int_0^1P_n(t)\Differential t
 =K\rho_1^n(1+o(1)),
 \qquad K>0,
 \label{eq:FM}
\end{equation}
which fixes the arithmetic-mean exponent \(\kappa_1\).

\subsection{Root mean square: an exact identity}

\begin{proposition}[Exact \(L^2\) norm]
\label{prop:L2}
For every \(n\ge1\),
\begin{equation}
 \int_0^1P_n(t)^2\Differential t=2^{-n}.
 \label{eq:L2}
\end{equation}
\end{proposition}

\begin{proof}
Use \(\sin^2u=(1-\cos2u)/2\):
\[
 P_n(t)^2=2^{-n}\prod_{j=0}^{n-1}
 (1-\cos(2\pi2^jt)).
\]
After expansion and product-to-sum, every nonconstant frequency is a signed
sum of distinct powers of two.  Uniqueness of binary representation prevents
such a sum from being zero.  Thus only the constant term survives integration.
\end{proof}
Consequently \(\rho_2=1/2\), \(\Lambda_2=2^{-1/2}\), and
\(\kappa_2=\log_2(2\pi)\).  This explains the exponent proposed in the
pre-wave draft: it is exactly the RMS exponent, not the pointwise or absolute
mean exponent.

\subsection{Uniform norm and the maximizing two-cycle}

Put \(u_j=\sin^2(\pi2^jt)\), so
\(u_{j+1}=4u_j(1-u_j)\), and define
\(V(u)=1+2u/3\).  The calibrated subaction underlying Gel'fond's bound\cite{gelfond} is the algebraic identity
\begin{equation}
 3V(u)-4uV(4u(1-u))
 =\frac{(2u+1)(4u-3)^2}{3}\ge0
 \label{eq:subaction-identity}
\end{equation}
implies
\begin{equation}
 |\sin(\pi2^jt)|
 \le\frac{\sqrt3}{2}
 \left(\frac{V(u_j)}{V(u_{j+1})}\right)^{1/2}.
 \label{eq:one-step-Gelfond}
\end{equation}
Telescoping gives
\begin{equation}
 P_n(t)\le\sqrt{\frac53}\left(\frac{\sqrt3}{2}\right)^n.
 \label{eq:Gelfond}
\end{equation}
At \(t=1/3\) or \(2/3\), every factor equals \(\sqrt3/2\).  Hence
\begin{equation}
 \Lambda_\infty=\frac{\sqrt3}{2},
 \qquad
 \kappa_\infty=\frac32+\frac{\log(\pi/\sqrt3)}{a}.
 \label{eq:Linfty}
\end{equation}
Combining \eqref{eq:Gelfond} with \eqref{eq:master-normalization} yields
\begin{equation}
 |\Phi(x)|\le C\Ecal_{\kappa_\infty}(x)
 \qquad(x\ge1),
 \label{eq:pointwise-upper}
\end{equation}
and equality in exponential scale on the fixed dyadic rays corresponding to
the period-two orbit, for example \(x=2^k(4/3)\).

\section{The natural arithmetic-mean gauge}
\label{sec:natural-gauge}

Let
\begin{equation}
 g_1(x):=\Ecal_{\kappa_1}(x).
 \label{eq:g1}
\end{equation}
By \eqref{eq:master-normalization} and the defining relation for \(\kappa_1\),
there is a continuous bounded mantissa weight \(W_1\), positive on \([0,1)\),
such that
\begin{equation}
 \frac{|\Phi(2^k(1+z))|}{g_1(2^k(1+z))}
 =\rho_1^{-(k+1)}W_1(z)P_{k+1}(z).
 \label{eq:natural-shell-density}
\end{equation}
RPF convergence implies that the finite measures on \([0,1]\)
\begin{equation}
 \Differential\eta_k(z):=
 \rho_1^{-(k+1)}W_1(z)P_{k+1}(z)\Differential z
 \label{eq:etak}
\end{equation}
converge weakly to a finite nonzero measure \(\eta\).  Define
\begin{equation}
 A_1:=\eta([0,1])
 =\lim_{k\to\infty}\int_0^1\frac{|\Phi(2^k(1+z))|}
 {g_1(2^k(1+z))}\Differential z.
 \label{eq:A1}
\end{equation}
A spectral computation gives \(A_1\approx0.09126612\), but only the exact
limit definition is needed.  Set
\begin{equation}
 g_0(x):=A_1g_1(x).
 \label{eq:g0}
\end{equation}

The measure \(\mu:=\eta/A_1\) is a probability measure.  It is atomless,
fully supported, and singular with respect to Lebesgue measure.  The
atomlessness and singularity are important enough to record explicitly.

\begin{proposition}[The limiting shell measure is continuous singular]
\label{prop:singular}
The probability measure \(\mu\) has no atoms, has full support on \([0,1]\),
and satisfies \(\mu\perp\Leb\).
\end{proposition}

\begin{proof}
The RPF eigenmeasure \(\nu\) for \(\LaplaceTransformSymbol_1^*\) obeys, at a point \(z\),
\(
 \rho_1\nu\{z\}=\tfrac12s(z)\nu\{\T z\}\).  Since
\(\rho_1>1/2\), a nonzero atom would force its forward masses to grow by a
factor at least \(2\rho_1>1\), contradicting finiteness.  Topological mixing
of doubling and positivity of the weight in the interior give full support.
Multiplication by the continuous positive factor \(W_1\) preserves these
properties.

For singularity, suppose the eigenmeasure had a nonzero absolutely continuous
part \(r\Differential z\).  The dual eigen-equation gives almost everywhere
\begin{equation}
 s(z)r(\T z)=\rho_1r(z).
 \label{eq:conformal-density}
\end{equation}
Thus
\begin{equation}
 r(\T^nz)=\frac{\rho_1^n}{P_n(z)}r(z).
 \label{eq:density-iterate}
\end{equation}
For Lebesgue-almost every \(z\), Birkhoff gives
\(P_n(z)=\exp[-n(a+o(1))]\), so on \(\{r>0\}\) the right side grows like
\((2\rho_1)^n\).  On the other hand, invariance of Lebesgue measure and
Markov's inequality give
\[
 \Leb\{r\circ\T^n>\EulerE^{\varepsilon n}\}
 =\Leb\{r>\EulerE^{\varepsilon n}\}
 \le\Norm{r}_1\EulerE^{-\varepsilon n}.
\]
Borel--Cantelli contradicts geometric growth when
\(0<\varepsilon<\log(2\rho_1)\).  The Lebesgue decomposition is preserved by
the dual operator, so the entire absolutely continuous part must vanish.
\end{proof}

\subsection{Dyadic endpoints, the phase profile, and boundedness}

For \(0\le z\le1\), let
\begin{equation}
 \mathscr F(z):=\mu([0,z]),
 \qquad
 \mathscr P(y):=\frac{1+\mathscr F(y-1)}{y},
 \quad 1\le y\le2.
 \label{eq:profile}
\end{equation}
The distribution function \(\mathscr F\) is continuous and singular, so
\(\mathscr P\) is continuous but not identically one.

\begin{theorem}[Natural-gauge Ces\`aro profile]
\label{thm:natural-profile}
For every fixed \(y\in[1,2]\),
\begin{equation}
 \frac1{2^ky-t_0}\int_{t_0}^{2^ky}
 \frac{|\Phi(x)|}{g_0(x)}\Differential x
 \longrightarrow\mathscr P(y).
 \label{eq:profile-limit}
\end{equation}
In particular, the limit is one at \(y=1,2\), and the ratio is eventually
bounded above and below by positive constants for unrestricted \(t\).
\end{theorem}

\begin{proof}
On a full shell, the change of variables \(x=2^k(1+z)\) and
\eqref{eq:natural-shell-density} show that the shell contribution divided by
\(2^k\) tends to one.  The sum of all completed shells below \(2^k\) is
therefore \(2^k+o(2^k)\).  The partial current shell up to \(z=y-1\) is
\(2^k(\mathscr F(z)+o(1))\), by weak convergence and atomlessness.  Division
by \(2^ky\) proves \eqref{eq:profile-limit}.  Continuity and positivity of
\(\mathscr P\) on the compact interval \([1,2]\) give the two-sided bound.
\end{proof}
The audit's numerical profile range, approximately
\([0.9249,1.1133]\), is consistent with this theorem but is not used in the
proof.

\subsection{What the singular profile really rules out}

\begin{theorem}[No exact limit for stabilizing shell shapes]
\label{thm:no-stabilizing}
Suppose \(g>0\) has the same rough shell scale as \(g_1\), and there are
constants \(c_k>0\) and a continuous \(u:[0,1]\to(0,\infty)\) such that
\begin{equation}
 \frac{g(2^k(1+z))}{c_kg_1(2^k(1+z))}
 \longrightarrow u(z)
 \quad\hbox{uniformly for }z\in[0,1].
 \label{eq:stabilizing}
\end{equation}
Then \Ceq cannot hold for \(g\).
\end{theorem}

\begin{proof}
If \Ceq held, subtraction of its asymptotics at \(2^k(1+z)\) and
\(2^k\) would force the normalized mass of every partial shell \([0,z]\) to
converge to \(z\).  By \eqref{eq:stabilizing} and RPF convergence, those
partial masses instead converge, after a scalar normalization, to
\(q\Differential\mu\), where \(q\) is continuous and strictly positive on \([0,1)\).
Uniform linearity would therefore imply \(q\Differential\mu=\Differential z\).  This is impossible
because \(\mu\) is singular and atomless.
\end{proof}

The class includes every Hardy-field gauge of the correct scale for which the
within-shell logarithmic derivative has a finite asymptotic slope, and also
such gauges multiplied by a continuous positive function of
\(\{\log_2x\}\).  Degenerate infinite slopes concentrate the shell mass at an
edge and fail the required linear partial-shell law even more directly.
This covers the usual logarithmico-exponential closed forms.  It does
\emph{not} cover every real-analytic decreasing function: analyticity permits
shell-dependent corrections whose complexity increases with \(k\).

\subsection{The logarithmic average repairs the natural gauge}

\begin{theorem}[Unrestricted logarithmic mean]
\label{thm:logmean}
There is a constant \(A_1^{\log}>0\) such that
\begin{equation}
 \frac1{\log(t/t_0)}\int_{t_0}^{t}
 \frac{|\Phi(x)|}{A_1^{\log}g_1(x)}\frac{\Differential x}{x}
 \longrightarrow1
 \qquad(t\to\infty)
 \label{eq:logmean}
\end{equation}
without any restriction on the dyadic mantissa of \(t\).
\end{theorem}

\begin{proof}
The logarithmic mass of the \(k\)-th shell converges by RPF to a positive
constant.  Summing \(K\) completed shells gives a quantity asymptotic to that
constant times \(K\), while the terminal partial shell is uniformly bounded.
Since \(\log t=Ka+\BigO(1)\), ordinary Ces\`aro summation proves the result.
The terminal shell has relative weight \(\BigO(1/\log t)\), so its singular
within-shell distribution disappears.
\end{proof}

\section{A literal all-time analytic gauge: shell-adaptive flattening}
\label{sec:adaptive}

This section closes the only remaining logical gap.  The natural gauge cannot
have an all-time arithmetic Ces\`aro limit because its normalized shell
measures tend to a singular measure.  But the original problem asks for
\emph{some} analytic decreasing gauge, not necessarily a gauge whose shell
shape stabilizes.  The deterministic baseline is so steep that a slowly
increasing amount of within-shell reweighting can be inserted without making
the total reciprocal gauge decrease.

Let
\begin{equation}
 D_k(z):=\frac{|\Phi(2^k(1+z))|}{g_0(2^k(1+z))},
 \qquad
 \Differential\mu_k(z):=D_k(z)\Differential z.
 \label{eq:muk}
\end{equation}
By construction and Theorem~\ref{thm:natural-profile},
\begin{equation}
 \mu_k\Longrightarrow\mu,
 \qquad \mu_k([0,1])\longrightarrow1,
 \label{eq:muk-conv}
\end{equation}
where \(\mu\) is atomless and fully supported.

\subsection{A flattening lemma for weakly convergent shell measures}

\begin{lemma}[Smooth asymptotic flattening]
\label{lem:flattening}
Let \(\mu_k\) be finite positive Borel measures on \([0,1]\) converging weakly
to an atomless, fully supported probability measure \(\mu\), with
\(\mu_k([0,1])\to1\).  Then there are functions
\(r_k\in C^\infty([0,1])\), \(r_k>0\), such that
\begin{enumerate}[label=\textup{(\alph*)}]
\item \(r_k\equiv1\) in neighborhoods of \(0\) and \(1\);
\item
\begin{equation}
 \sup_{0\le z\le1}
 \left|\int_0^zr_k\Differential\mu_k-z\right|\longrightarrow0;
 \label{eq:flattening-cdf}
\end{equation}
\item
\begin{equation}
 \Norm{\log r_k}_\infty=o(k),
 \qquad
 \Norm{(\log r_k)'}_\infty=o(k).
 \label{eq:flattening-slopes}
\end{equation}
\end{enumerate}
\end{lemma}

\begin{proof}
We give the diagonal construction because it is the mechanism that defeats the
stabilizing-class obstruction.

For each integer \(j\ge1\), choose a finite partition
\(0=x_{j,0}<x_{j,1}<\cdots<x_{j,N_j}=1\) with mesh at most
\(\varepsilon_j\downarrow0\), and with all boundary points non-atoms of
\(\mu\) (automatic here).  Full support gives
\(\mu(I_{j,i})>0\) for every partition interval
\(I_{j,i}=[x_{j,i-1},x_{j,i}]\).  Define limiting cell multipliers
\[
 c_{j,i}:=\frac{|I_{j,i}|}{\mu(I_{j,i})}.
\]
Choose a finite union \(U_j\) of very small intervals around all partition
boundaries, including \(0,1\), so that both \(\Leb(U_j)\) and \(\mu(U_j)\)
are as small as needed below and \(\mu(\partial U_j)=0\).

For sufficiently large \(k\), all cell masses \(\mu_k(I_{j,i})\) are close
to \(\mu(I_{j,i})\), and \(\mu_k(U_j)\) is close to \(\mu(U_j)\).  Put
\[
 c_{k,j,i}:=\frac{|I_{j,i}|}{\mu_k(I_{j,i})}.
\]
On the portion of each cell outside \(U_j\), set
\(\ell_{k,j}=\log c_{k,j,i}\).  Join the finitely many constants smoothly
inside \(U_j\), and arrange \(\ell_{k,j}=0\) near \(0\) and \(1\).
For fixed \(j\), the values \(c_{k,j,i}\) stay in a compact subinterval of
\((0,\infty)\), so the interpolation can be chosen with uniform bounds
\[
 \Norm{\ell_{k,j}}_\infty\le B_j,
 \qquad
 \Norm{\ell_{k,j}'}_\infty\le C_j
 \quad(k\ge K_j^{(0)}).
\]
Let \(r_{k,j}=\EulerE^{\ell_{k,j}}\).  Away from \(U_j\), the weighted mass of a
whole cell is exactly its Lebesgue length.  The discrepancy is confined to
\(U_j\), where the multiplier is uniformly bounded.  Taking \(U_j\) small
enough therefore gives, for all sufficiently large \(k\),
\[
 \sup_z\left|\int_0^zr_{k,j}\Differential\mu_k-z\right|
 \le 3\varepsilon_j.
\]
The mesh bounds the possible error in the one partially traversed cell.

Choose increasing integers \(K_j\) so large that all preceding conclusions
hold for \(k\ge K_j\), and additionally
\[
 \frac{B_j}{K_j}\le\varepsilon_j,
 \qquad
 \frac{C_j}{K_j}\le\varepsilon_j.
\]
For \(K_j\le k<K_{j+1}\), set \(r_k=r_{k,j}\).  Then
\(j=j(k)\to\infty\), so the cumulative discrepancy tends to zero and both
ratios in \eqref{eq:flattening-slopes} tend to zero.  The endpoint condition
was built into every interpolation.
\end{proof}

\subsection{A positive smooth decreasing exact gauge}

Define a global multiplier \(R\) on all sufficiently large \(x\) by
\begin{equation}
 R(2^k(1+z)):=r_k(z),
 \qquad 0\le z\le1.
 \label{eq:R-global}
\end{equation}
Because every \(r_k\) equals one in neighborhoods of both endpoints, these
pieces agree as the constant function one in neighborhoods of every dyadic
boundary.  Thus \(R\) is globally \(C^\infty\).

Let
\begin{equation}
 h_0(x):=\frac1{g_0(x)},
 \qquad
 h(x):=h_0(x)R(x),
 \qquad
 g_{\infty}(x):=\frac1{h(x)}.
 \label{eq:smooth-adaptive}
\end{equation}
In shell coordinates,
\begin{equation}
 \frac{\partial}{\partial z}
 \log h_0(2^k(1+z))
 =\frac{k+\kappa_1+a^{-1}\log(1+z)}{1+z}.
 \label{eq:baseline-slope}
\end{equation}
The right side is at least \(k/3\) for all large \(k\), uniformly in
\(z\in[0,1]\).  By \eqref{eq:flattening-slopes},
\(\Norm{(\log r_k)'}_\infty=o(k)\); hence \(h\) is eventually strictly
increasing and \(g_\infty\) is eventually strictly decreasing.  Also
\begin{equation}
 \log g_\infty(x)=\log g_0(x)+o(\log x).
 \label{eq:smooth-rough}
\end{equation}

\begin{proposition}[Exact arithmetic mean for the smooth adaptive gauge]
\label{prop:smooth-adaptive}
The gauge \(g_\infty\) satisfies
\begin{equation}
 \frac1{t-t_0}\int_{t_0}^{t}
 \frac{|\Phi(x)|}{g_\infty(x)}\Differential x\longrightarrow1
 \qquad(t\to\infty).
 \label{eq:smooth-exact}
\end{equation}
\end{proposition}

\begin{proof}
Let
\[
 C_k(z):=\int_0^zD_k(u)r_k(u)\Differential u.
\]
Lemma~\ref{lem:flattening} gives
\(\sup_z|C_k(z)-z|\to0\).  If
\(t=2^K(1+z)\), then, up to a fixed initial contribution,
\[
 \int_{t_0}^{t}\frac{|\Phi(x)|}{g_\infty(x)}\Differential x
 =\sum_{k<K}2^kC_k(1)+2^KC_K(z).
\]
Replacing \(C_k(1)\) by \(1\) and \(C_K(z)\) by \(z\) creates an
\(o(2^K)\) error: this is the elementary Toeplitz fact that
\(\sum_{k<K}2^k\varepsilon_k=o(2^K)\) whenever \(\varepsilon_k\to0\).
The main term is
\(\sum_{k<K}2^k+2^Kz=2^K(1+z)+\BigO(1)=t+\BigO(1)\).
This proves \eqref{eq:smooth-exact}.
\end{proof}

\subsection{Real-analytic upgrade by Carleman--Hoischen approximation}

We use the following classical approximation theorem in its \(C^1\) form:
for every real-valued \(F\in C^1(\RealNumbers)\) and every positive continuous
\(\epsilon\), there is an entire function \(G\), real on \(\RealNumbers\), such that
\begin{equation}
 |G(x)-F(x)|<\epsilon(x),
 \qquad
 |G'(x)-F'(x)|<\epsilon(x)
 \quad(x\in\RealNumbers).
 \label{eq:Hoischen}
\end{equation}
This is Hoischen's derivative version of Carleman approximation; a modern
proof is due to Gauthier and Kienzle\cite{gauthier-kienzle}.

Extend \(h\) from \eqref{eq:smooth-adaptive} to a positive strictly increasing
\(C^1\) function on all of \(\RealNumbers\).  Choose a continuous
\(\delta(x)>0\) with \(\delta(x)\to0\) as \(x\to+\infty\), and apply
\eqref{eq:Hoischen} with
\begin{equation}
 \epsilon(x)<\frac12\min\{h(x),h'(x),\delta(x)h(x)\}.
 \label{eq:Hoischen-error}
\end{equation}
Then the real entire approximant \(\widetilde h\) satisfies
\(\widetilde h>0\), \(\widetilde h'>0\) on \(\RealNumbers\), and
\(\widetilde h/h\to1\) as \(x\to+\infty\).  Define
\begin{equation}
 g_{\mathrm{ad}}(x):=\frac1{\widetilde h(x)},
 \qquad x>0.
 \label{eq:g-ad}
\end{equation}
It is positive, real-analytic, and strictly decreasing.

\begin{theorem}[Literal analytic solution of the original limit problem]
\label{thm:analytic-adaptive}
The gauge \(g_{\mathrm{ad}}\) in \eqref{eq:g-ad} satisfies \Ceq and
\eqref{eq:adaptive-rough}.
\end{theorem}

\begin{proof}
The rough asymptotic follows from
\(\widetilde h/h\to1\), \eqref{eq:smooth-rough}, and the definition of
\(g_0\).  For the integral, given \(\varepsilon>0\), choose \(X\) so that
\(|\widetilde h/h-1|<\varepsilon\) for \(x\ge X\).  Then
\[
 \left|\int_X^t|\Phi|(\widetilde h-h)\Differential x\right|
 \le\varepsilon\int_X^t|\Phi|h\Differential x
 =\varepsilon(t+o(t))
\]
by Proposition~\ref{prop:smooth-adaptive}.  The initial interval contributes
\(\BigO(1)\).  Letting \(\varepsilon\downarrow0\) proves the same limit for
\(\widetilde h=1/g_{\mathrm{ad}}\).
\end{proof}

\begin{remark}[Why this does not contradict the natural obstruction]
The multiplier \(r_k\) corrects only finitely many cells at level \(k\), but
the number of cells tends to infinity.  Its logarithm and logarithmic slope
are \(o(k)\), so the total gauge remains decreasing and retains the
\(\kappa_1\) rough exponent.  Nevertheless its shell shape never converges.
Theorem~\ref{thm:no-stabilizing} therefore does not apply.  This also explains
why no fixed power of \(x\), fixed power of \(\log x\), or continuous
log-periodic multiplier can produce the exact arithmetic limit.
\end{remark}

\begin{remark}[On the optional higher-derivative request]
The natural gauge \(g_0\) has alternating derivative signs eventually for
each fixed order.  The adaptive analytic construction above guarantees
positivity and strict decrease, which are the mandatory regularity conditions
in the original wording.  It does not claim one common tail on which every
higher derivative is monotone.  Preserving a prescribed finite number of
higher derivative signs is possible by using the corresponding higher-order
Carleman--Hoischen theorem and slowing the diagonal construction further.
\end{remark}

\section{The sharp pointwise and annular envelopes}
\label{sec:annular}

The pointwise upper scale \(\Ecal_{\kappa_\infty}(x)\) is sharp along fixed
maximizing rays, but the largest value in the moving annulus
\([2^k,2^{k+1}]\) is a subtler boundary-layer problem.  A fixed phase such as
\(y=4/3\) incurs an exponential loss relative to the left endpoint scale.
The true maximizer approaches the zero at \(2^k\), spends a slowly growing
number of doubling steps in the small-angle regime, and then shadows the
maximizing cycle \(\{1/3,2/3\}\).

Put
\begin{equation}
 M_k:=\max_{2^k\le x\le2^{k+1}}|\Phi(x)|,
 \qquad n=k+1,
 \qquad
 \beta:=\log\frac{\sqrt3}{2},
 \label{eq:Mk-beta}
\end{equation}
and
\begin{equation}
 c:=\frac{a}{\pi}\sqrt{\frac32}.
 \label{eq:c}
\end{equation}

\subsection{Boundary-layer reduction}

Writing \(x=2^k(1+\theta)\) in \eqref{eq:factorization} gives
\begin{equation}
 |\Phi(2^k(1+\theta))|
 =2^{-k(k+1)/2}\pi^{-n}
 H(1+\theta)(1+\theta)^{-n}P_n(\theta).
 \label{eq:annulus-exact}
\end{equation}
The symmetry \(P_n(\theta)=P_n(1-\theta)\) reduces the supremum, up to a
constant independent of \(n\), to \(0\le\theta\le1/2\).  Indeed, if
\(1/2\le\theta\le1\), then \(H(2-\theta)\) is bounded below while
\(H(1+\theta)\) is bounded above, and
\((2-\theta)/(1+\theta)\le1\).  Consequently the value at \(\theta\) is at
most a fixed multiple of the value at \(1-\theta\).

For \(0<\theta\le1/2\), write
\begin{equation}
 \theta=2^{-r}s,
 \qquad r\in\IntegerNumbers_{\ge0},
 \qquad s\in I:=[1/4,1/2].
 \label{eq:transition}
\end{equation}
If \(r<n\), then
\begin{equation}
 P_n(2^{-r}s)
 =\left[\prod_{h=1}^{r}\sin\!\left(\frac{\pi s}{2^h}\right)\right]
 P_{n-r}(s).
 \label{eq:transition-product}
\end{equation}
Uniformly for \(s\in I\), the initial block satisfies
\begin{equation}
 \log\prod_{h=1}^{r}\sin\!\left(\frac{\pi s}{2^h}\right)
 =-\frac a2r^2+r\left(\log(\pi s)-\frac a2\right)+\BigO(1).
 \label{eq:small-block}
\end{equation}
The Gelfond inequality gives the upper tail
\(P_m(s)\le C\EulerE^{m\beta}\).  Conversely, the \(q\)-fold preimages of
\(1/3\) and \(2/3\) form an \(O(2^{-q})\)-net.  For example, if
\(s_q=(j+1/3)/2^q\), then \(\T^qs_q=1/3\), and for \(0\le \ell<q\)
\[
 \operatorname{dist}(\T^\ell s_q,\IntegerNumbers)
 \ge \frac1{3\,2^{q-\ell}}.
\]
Since \(|\sin\pi u|\ge2\operatorname{dist}(u,\IntegerNumbers)\), the first \(q\)
factors are bounded below by
\[
 \prod_{\ell=0}^{q-1}|\sin(\pi2^\ell s_q)|
 \ge \left(\frac23\right)^q2^{-q(q+1)/2}.
\]
All later factors lie exactly on the maximizing two-cycle.  Hence a nearby
preimage can be chosen with
\begin{equation}
 P_m(s_q)\ge\exp(m\beta-Cq^2),
 \qquad m\ge q.
 \label{eq:dense-preimages}
\end{equation}
At the relevant transition indices \(r=\log_2 n+O(\log\log n)\), the finite
objective is \(O(\log n)\)-Lipschitz in \(s\).  Taking
\(q=\lceil2\log_2\log n\rceil\) therefore makes the net-discretization loss
\(o(1)\), while the initial-orbit loss is \(O(q^2)\).  This is the final
\(O((\log\log n)^2)\) error term.

The finite objective is
\begin{equation}
 \DecayOptimizationObjective{n}(r,s)
 =-\frac a2r^2
 +r\left(\log(\pi s)-\frac a2-\beta\right)
 -n\log(1+s2^{-r}).
 \label{eq:F-objective}
\end{equation}
The preceding estimates give
\begin{equation}
 \log\sup_{0\le\theta\le1}
 H(1+\theta)(1+\theta)^{-n}P_n(\theta)
 =n\beta+\sup_{r,s}\DecayOptimizationObjective{n}(r,s)
 +\BigO((\log\log n)^2).
 \label{eq:Fn-reduction}
\end{equation}

\subsection{Lambert optimization}

Set
\begin{equation}
 C_0:=\log\pi-\frac a2-\beta
 =\log\!\left(\pi\sqrt{\frac23}\right).
 \label{eq:C0}
\end{equation}
Replacing the final logarithm in \eqref{eq:F-objective} by its linear term
changes the supremum by \(o(1)\).  Put
\(A=C_0+\log s\) and \(z=ar-A\).  As \(s\) ranges over \([1/4,1/2]\),
\(A\) ranges over an interval of length exactly \(a\); as \(r\) runs through
the integers, the corresponding \(z\)-intervals tile the real line.  Up to a
bounded term, the objective becomes
\begin{equation}
 -\frac{z^2}{2a}-n\EulerE^{-C_0}\EulerE^{-z}.
 \label{eq:z-objective}
\end{equation}
Its critical equation is
\begin{equation}
 z\EulerE^z=an\EulerE^{-C_0}=cn.
 \label{eq:Lambert-eq}
\end{equation}
Thus \(z=W(cn)\), with the principal Lambert function\cite{corless}, and at the maximum
\begin{equation}
 \sup_{r,s}\DecayOptimizationObjective{n}(r,s)
 =-\frac{W(cn)^2+2W(cn)}{2a}+\BigO(1).
 \label{eq:Lambert-sup}
\end{equation}

\begin{theorem}[Sharp dyadic-annulus maximum]
\label{thm:annulus}
As \(k\to\infty\), with \(n=k+1\),
\begin{equation}
 \boxed{
 \log M_k
 =-\frac a2k(k+1)-n\log\pi+n\beta
 -\frac{W(cn)^2+2W(cn)}{2a}
 +\BigO((\log\log k)^2).}
 \label{eq:annulus-main}
\end{equation}
\end{theorem}

Expanding \(W(cn)=\log n-\log\log n+\log c+\BigO(\log\log n/\log n)\)
gives
\begin{align}
 \log M_k={}&-\frac a2k^2
 -\left(\frac a2+\log\pi-\beta\right)k
 -\frac{(\log k)^2}{2a}
 +\frac{\log k\log\log k}{a}\notag\\
 &-\frac{1+\log c}{a}\log k
 +\BigO((\log\log k)^2).
 \label{eq:annulus-expanded}
\end{align}
In physical scale \(R=2^k\), the extra correction begins with
\begin{equation}
 -\frac{(\log\log R)^2}{2\log2}
 +\frac{\log\log R\,\log\log\log R}{\log2}
 +\BigO(\log\log R).
 \label{eq:second-lognormal}
\end{equation}
No fixed factor \((\log R)^{-p}\) can represent this boundary layer.  The
near-maximizing phase satisfies
\begin{equation}
 \theta_k\asymp\frac{W(c(k+1))}{k+1},
 \qquad
 x_k=2^k\left(1+\Theta\!\left(\frac{\log k}{k}\right)\right).
 \label{eq:max-location}
\end{equation}

\section{Exceptional local peaks and two-adic valleys}
\label{sec:valleys}

Let \(E_N\) denote the unique maximum of \(|\Phi|\) on \((N,N+1)\).
Fix an odd positive integer \(m\), put \(N=m2^k\), and write
\(x=N+t\), \(0\le t\le1\).  The first \(k+1\) factors yield uniformly
\begin{equation}
 |\Phi(m2^k+t)|
 =(m2^k)^{-(k+1)}|H(m)|\,t^{k+1}\Phi(t)(1+o(1)).
 \label{eq:near-m2k}
\end{equation}
Near \(t=1\),
\begin{equation}
 \Phi(t)=H(1)(1-t)+\BigO((1-t)^2).
 \label{eq:Phi-near-one}
\end{equation}
The elementary boundary maximum
\(\max_{0\le t\le1}t^p(1-t)\sim1/(\EulerE p)\) now gives the exact valley law.

\begin{theorem}[Peaks after a fixed odd multiple of a power of two]
\label{thm:twoadic}
Let \(c_{m2^k}=m2^k+t_k\) be the unique critical point in
\((m2^k,m2^k+1)\).  Then
\begin{equation}
 (k+1)(1-t_k)\longrightarrow1,
 \label{eq:twoadic-location}
\end{equation}
and
\begin{equation}
 \boxed{
 E_{m2^k}\sim
 \frac{H(1)|H(m)|}{\EulerE(k+1)}(m2^k)^{-(k+1)}.}
 \label{eq:twoadic-peak}
\end{equation}
In particular,
\begin{equation}
 E_{2^k}\sim\frac{H(1)^2}{\EulerE(k+1)}2^{-k(k+1)}.
 \label{eq:power-two-peak}
\end{equation}
\end{theorem}

The leading quadratic coefficient in \(\log E_{2^k}\) is \(-\log2\),
twice the magnitude of the annular maximum's coefficient
\(-\tfrac12\log2\).  Thus even the sequence of local maxima has no single
two-sided log-normal--power envelope.

\section{Implications, limitations, and formalization targets}

\subsection{What is now completely settled}

The following conclusions are rigorous, with their external analytic inputs
explicitly identified.
\begin{itemize}
\item The exact factorization, the \(q=0\), \(q=2\), and \(q=\infty\) rows,
the variance correction, the Gelfond inequality, the Lambert optimization,
and the two-adic valley asymptotic are elementary consequences of the product
plus standard real analysis.
\item The \(q=1\) exponent uses the Fouvry--Mauduit asymptotic
\eqref{eq:FM}.  The shell profile uses the standard RPF spectral theorem for
\(\LaplaceTransformSymbol_1\).
\item The analytic adaptive gauge additionally uses the \(C^1\)
Carleman--Hoischen approximation theorem.  Its shell-flattening construction
is otherwise elementary once weak convergence \(\mu_k\Rightarrow\mu\) is
known.
\item The literal original existence question \Ceq and the weaker
question \Beq are both affirmative.  What fails is not existence, but
existence inside the asymptotically stationary closed-form class.
\end{itemize}

\subsection{What should not be claimed}

There is no single pointwise equivalent, no universal exponent independent of
the chosen norm, and no elementary justification for replacing the lacunary
numerator by a typical mean when studying maxima.  The high-precision decimals
for \(\rho_1\), \(A_1\), and the phase profile are numerical, although the
operators and limiting constants are exact.  The analytic adaptive gauge is an
existence construction, not a short elementary formula.  Finally, the
higher-derivative monotonicity requested only optionally in the original
question is not part of Theorem~\ref{thm:analytic-adaptive}.

\subsection{Lean-friendly decomposition}

A formal development can be divided into the following blocks.
\begin{enumerate}
\item Locally uniform convergence, refinement, zeros, and the exact shell
factorization.
\item Strict log-concavity and uniqueness of every interzero maximum.
\item Exact \(L^2\) identity and the calibrated polynomial inequality
\eqref{eq:subaction-identity}.
\item Fourier covariance calculation proving \(\sigma^2=\pi^2/4\).
\item Rational preimages and the finite Lambert objective.
\item Boundary maximum lemma and two-adic valleys.
\item Imported abstractions for Birkhoff, martingale CLT/LIL, Fouvry--Mauduit,
and RPF.
\item The shell-flattening diagonal lemma, which is measure-theoretic and
independent of the special product after \(\mu_k\Rightarrow\mu\) is supplied.
The final analytic upgrade may remain as an external complex-approximation
theorem.
\end{enumerate}

\section{Conclusion}

The Fourier transform of Rvachev's up-function does have an exact dominant
decay scale:
\[
 \exp\!\left(-\frac{(\log x)^2}{2\log2}\right).
\]
Everything beyond that universal factor records how the lacunary sine cocycle
is sampled.  A typical orbit, an \(L^1\) average, an \(L^2\) average, and an
extremal orbit yield four distinct powers.  The largest peak in a moving
dyadic annulus is more delicate still, acquiring a second log-normal layer in
\(\log\log x\).  The correct fluctuation variance is \(\pi^2/4\), resolving
the only numerical disagreement among the first-wave analyses and exposing a
Parseval slip in the cited 2017 paper.

For the exact question originally asked, the distinction between a
\emph{natural} gauge and an arbitrary analytic gauge is decisive.  The natural
transfer-operator gauge is explicit up to one spectral constant, solves the
bounded problem, and displays a singular dyadic phase profile.  That profile
rules out every asymptotically stationary closed form.  It does not rule out
analyticity itself.  By allowing the gauge to adapt on successively finer
shell partitions---slowly enough that the universal log-square slope still
dominates---one obtains a positive analytic decreasing gauge with the full
all-time Ces\`aro limit.  This is the final literal resolution: the desired
analytic function exists, has the arithmetic-mean exponent \(\kappa_1\), but
cannot be represented by a fixed elementary correction to the log-normal
baseline.

\appendix

\section{Constants and conversions}

For convenience,
\begin{align*}
 a&=\log2=0.6931471805599453094\ldots,\\
 \beta&=\log(\sqrt3/2)=-0.1438410362258904637\ldots,\\
 c&=\frac{a}{\pi}\sqrt{\frac32}
   =0.2702223197333653409\ldots,\\
 \frac a2+\log\pi-\beta
  &=1.6351445123552632926\ldots,\\
 \kappa_0&=3.1514961294723187980\ldots,\\
 \kappa_1&=2.7480700148713340157\ldots,\\
 \kappa_2&=2.6514961294723187980\ldots,\\
 \kappa_\infty&=2.3590148791117407073\ldots.
\end{align*}
The LIL conversions are
\[
 \frac\pi2\sqrt{2n\log\log n}
 =\frac\pi{\sqrt2}\sqrt{n\log\log n},
\]
and, because \(n\sim(\log x)/\log2\), fixed-ray LIL fluctuations in the
frequency variable have order
\[
 \sqrt{\log x\,\log\log\log x}.
\]

\section{Source map}

The repository sources audited in this article are listed below so that every
claim can be traced to its originating draft.  The present article changes two
points relative to the comparative audit: it retains the factor-\(\sqrt2\)
LIL correction, and it replaces the unrestricted nonexistence claim by the
shell-adaptive analytic existence theorem of Section~\ref{sec:adaptive}.

\begin{thebibliography}{99}

\bibitem{arias1982}
J.~Arias de Reyna,
\emph{An infinitely differentiable function with compact support: definition and properties},
English translation of the 1982 article, arXiv:1702.05442,
\url{https://arxiv.org/abs/1702.05442}.

\bibitem{ahl2017}
C.~Aistleitner, R.~Hofer, and G.~Larcher,
\emph{On evil Kronecker sequences and lacunary trigonometric products},
Ann. Inst. Fourier (Grenoble) \textbf{67} (2017), no.~2, 637--687;
preprint arXiv:1502.06738,
\url{https://arxiv.org/abs/1502.06738}.

\bibitem{fouvrymauduit}
\'E.~Fouvry and C.~Mauduit,
\emph{Sommes des chiffres et nombres presque premiers},
Math. Ann. \textbf{305} (1996), 571--599.

\bibitem{baladi}
V.~Baladi,
\emph{Positive Transfer Operators and Decay of Correlations},
World Scientific, 2000.

\bibitem{hallheyde}
P.~Hall and C.~C. Heyde,
\emph{Martingale Limit Theory and Its Application},
Academic Press, 1980.

\bibitem{gelfond}
A.~O. Gel'fond,
\emph{Sur les nombres qui ont des propri\'et\'es additives et multiplicatives donn\'ees},
Acta Arith. \textbf{13} (1967/68), 259--265.

\bibitem{corless}
R.~M. Corless, G.~H. Gonnet, D.~E. Hare, D.~J. Jeffrey, and D.~E. Knuth,
\emph{On the Lambert \(W\) function},
Adv. Comput. Math. \textbf{5} (1996), 329--359.

\bibitem{gauthier-kienzle}
P.~M. Gauthier and J.~Kienzle,
\emph{Approximation of a function and its derivatives by entire functions},
Canad. Math. Bull. \textbf{59} (2016), no.~1, 87--94,
\url{https://doi.org/10.4153/CMB-2015-060-5}.

\bibitem{mse-question}
V.~Reshetnikov,
\emph{Asymptotic decay rate of an infinite product of sinc functions},
Mathematics Stack Exchange question 2745497 (2018),
\url{https://math.stackexchange.com/questions/2745497/asymptotic-decay-rate-of-an-infinite-product-of-sinc-functions}.

\bibitem{question-source}
V.~Reshetnikov,
\emph{Asymptotic decay rate of an infinite product of sinc functions},
repository transcription of the original question,
\url{https://github.com/VladimirReshetnikov/ProveIt/blob/main/Analysis/FabiusFunction/docs/semi-formalized-research-frontiers/drafts/rvachev_up_fourier_decay/asymptotic-decay-rate-of-an-infinite-product-of-sinc-functions/asymptotic-decay-rate-of-an-infinite-product-of-sinc-functions.tex}.

\bibitem{wave1-1}
\emph{Decay of the Fourier Transform of Rvachev's Up-Function: Dyadic products,
ergodic fluctuations, and sharp annular envelopes}, first-wave report 1,
\url{https://github.com/VladimirReshetnikov/ProveIt/blob/main/Analysis/FabiusFunction/docs/semi-formalized-research-frontiers/drafts/rvachev_up_fourier_decay/rvachev_up_fourier_decay-1/rvachev_up_fourier_decay.tex}.

\bibitem{wave1-2}
\emph{Sharp Real-Axis Decay of the Fourier Transform of Rvachev's Up-Function},
first-wave report 2,
\url{https://github.com/VladimirReshetnikov/ProveIt/blob/main/Analysis/FabiusFunction/docs/semi-formalized-research-frontiers/drafts/rvachev_up_fourier_decay/Rvachev_Up_Fourier_Decay-2/Rvachev_Up_Fourier_Decay.tex}.

\bibitem{wave1-3}
\emph{Sharp Decay Laws for the Fourier Transform of Rvachev's Up-Function},
first-wave report 3,
\url{https://github.com/VladimirReshetnikov/ProveIt/blob/main/Analysis/FabiusFunction/docs/semi-formalized-research-frontiers/drafts/rvachev_up_fourier_decay/Rvachev_Up_Fourier_Decay-3/Fourier_decay_of_Rvachev_up.tex}.

\bibitem{wave1-4}
\emph{Decay of the Fourier Transform of Rvachev's Up-Function: Dyadic
renormalization, transfer-operator norms, sharp envelopes, and exceptional valleys},
first-wave report 4,
\url{https://github.com/VladimirReshetnikov/ProveIt/blob/main/Analysis/FabiusFunction/docs/semi-formalized-research-frontiers/drafts/rvachev_up_fourier_decay/rvachev_up_fourier_decay-4/rvachev_up_fourier_decay.tex}.

\bibitem{audit}
\emph{Rvachev Up Fourier Decay: Comparative Audit}, second-wave report,
\url{https://github.com/VladimirReshetnikov/ProveIt/blob/main/Analysis/FabiusFunction/docs/semi-formalized-research-frontiers/drafts/rvachev_up_fourier_decay/Rvachev_Up_Fourier_Decay_Comparative_Audit/Rvachev_Up_Fourier_Decay_Comparative_Audit.tex}.

\bibitem{ahl-source}
C.~Aistleitner, R.~Hofer, and G.~Larcher,
LaTeX source of arXiv:1502.06738v2 as stored in \texttt{ProveIt},
\url{https://github.com/VladimirReshetnikov/ProveIt/blob/main/Analysis/FabiusFunction/docs/papers/arXiv-1502.06738v2/Lacunary_products.tex}.

\bibitem{wikipedia}
\emph{Fabius function}, Wikipedia,
\url{https://en.wikipedia.org/wiki/Fabius_function}.

\end{thebibliography}

\end{document}
